% ==============================================================================
% Section VIII: FPGA Implementation Results, Performance Analysis & State-of-the-Art Comparison
% ==============================================================================
\section{FPGA Implementation Results, Performance Analysis \& State-of-the-Art Comparison}
\label{sec:results}

This section presents the physical implementation methodology, post-routing resource utilization, static timing closure, execution latency benchmarks, and comprehensive state-of-the-art comparison of the proposed unified \mbox{ML-KEM} hardware accelerator.

\subsection{Implementation Platform \& Experimental Methodology}
\label{sec:results_platform}

The complete top-level \mbox{ML-KEM} accelerator core (\texttt{mlkem\_keygen\_top}) and all constituent arithmetic, memory, and cryptographic submodules were implemented using the AMD Vivado Design Suite 2025.2 (Build 6299465, 64-bit Linux). The target physical silicon device is the \textbf{AMD Kintex UltraScale \texttt{xcku035-ffva1156-3-e}} FPGA (speed grade -3, extended temperature grade), fabricated on TSMC's 20nm planar high-$k$ metal gate (HKMG) process.

Implementation was performed under rigorous out-of-context (OOC) design constraints with physical logic synthesis, default placement directives, and high-effort physical optimization (\texttt{phys\_opt\_design -directive AggressiveExplore}). The primary synthesis clock constraint was set to an aggressive target period:
\begin{equation}
\label{eq:clock_constraint}
T_{\text{target}} = 2.222\text{ ns} \quad (F_{\text{target}} = 450.045\text{ MHz}).
\end{equation}

Functional verification was executed using the Vivado Simulator (\texttt{xsim}) across the complete parameter spectrum. Bit-exact compliance with the official NIST FIPS~203 standard~\cite{fips203} was verified against 100\% of golden Known Answer Test (KAT) vectors from the Cryptographic Algorithm Validation Program (CAVP) test suites for ML-KEM-512 ($k = 2$), ML-KEM-768 ($k = 3$), and ML-KEM-1024 ($k = 4$). In addition, randomized stress testing, corner-case injection, non-canonical public key coefficient detection, and constant-time implicit rejection validation were conducted across 495 complete decapsulation executions.

\subsection{Comprehensive Resource Utilization Breakdown}
\label{sec:results_utilization}

\subsubsection{Device-Level Physical Utilization}
Table~\ref{tab:device_utilization} summarizes the exact post-routing physical resource consumption extracted from the routed implementation report (\texttt{util\_impl.rpt}) on the AMD Kintex UltraScale device.

\begin{table}[htbp]
\caption{Comprehensive Post-Implementation Resource Utilization Summary on AMD Kintex UltraScale (\texttt{xcku035-ffva1156-3-e})}
\label{tab:device_utilization}
\centering
\scriptsize
\setlength{\tabcolsep}{2.0pt}
\begin{tabular}{lrrr}
\toprule
\textbf{Resource Primitive} & \textbf{Used Count} & \textbf{Available} & \textbf{Util. (\%)} \\
\midrule
\textbf{Total CLB LUTs} & \textbf{29,063} & \textbf{203,128} & \textbf{14.31\%} \\
--- \textit{Logic LUTs} & 23,484 & 203,128 & 11.56\% \\
--- \textit{Memory LUTs (Total)} & 5,579 & 112,800 & 4.95\% \\
\quad --- \textit{Distributed RAM (LUTRAM)} & 5,324 & 112,800 & 4.72\% \\
\quad --- \textit{Shift Register (SRL)} & 255 & 112,800 & 0.23\% \\
\textbf{CLB Flip-Flops (Registers)} & \textbf{17,666} & \textbf{406,256} & \textbf{4.35\%} \\
\textbf{Occupied CLB Slices} & \textbf{5,453} & \textbf{30,300} & \textbf{18.00\%} \\
--- \textit{CLBL Slices (Logic)} & 2,617 & --- & --- \\
--- \textit{CLBM Slices (Distributed Mem.)} & 2,836 & --- & --- \\
\textbf{Block RAM Tiles (RAMB36/18/FIFO)} & \textbf{0} & \textbf{540} & \textbf{0.00\% (0 BRAM)} \\
\textbf{DSP48E2 Slices} & \textbf{8} & \textbf{1,700} & \textbf{0.47\% (8 DSPs)} \\
Dedicated Carry (\texttt{CARRY8}) & 519 & 30,300 & 1.71\% \\
Multiplexers (\texttt{F7MUX/F8MUX}) & 560 & 181,800 & 0.31\% \\
Unique Control Sets & 339 & 60,600 & 0.56\% \\
\bottomrule
\end{tabular}
\end{table}

\subsubsection{Hierarchical Submodule Resource Breakdown}
Table~\ref{tab:submodule_utilization} details the post-routing physical resource consumption apportioned across the primary architectural submodules of the accelerator.

\begin{table}[htbp]
\caption{Post-Implementation Resource Breakdown by Sub-Module}
\label{tab:submodule_utilization}
\centering
\scriptsize
\setlength{\tabcolsep}{2.5pt}
\begin{tabular}{lrrrr}
\toprule
\textbf{Sub-Module / Functional Block} & \textbf{Logic} & \textbf{Mem.} & \textbf{Reg.} & \textbf{DSP} \\
\midrule
NTT/INTT Engine (\texttt{mlkem\_ntt\_intt}) & 3,842 & 96 & 2,854 & 0 \\
Pointwise Multiplier (\texttt{mlkem\_mul\_stream}) & 1,870 & 184 & 1,412 & 8 \\
Modular Field ALUs (\texttt{field\_add/sub}) & 130 & 0 & 78 & 0 \\
Cryptographic Sponge (\texttt{keccak\_f1600}) & 4,746 & 0 & 1,609 & 0 \\
Hash Sampler \& CBD (\texttt{mlkem\_hash\_sampler}) & 6,204 & 0 & 3,102 & 0 \\
Polynomial Codecs Subsystem & 602 & 0 & 692 & 0 \\
\quad --- Compressor (\texttt{mlkem\_poly\_compress}) & 142 & 0 & 148 & 0 \\
\quad --- Decompressor (\texttt{mlkem\_poly\_decompress}) & 88 & 0 & 82 & 0 \\
\quad --- Serialization Encoder (\texttt{poly\_encode}) & 178 & 0 & 234 & 0 \\
\quad --- Deserialization Decoder (\texttt{poly\_decode}) & 194 & 0 & 228 & 0 \\
Top Control \& Memory (\texttt{mlkem\_keygen\_top}) & 6,090 & 5,299 & 7,919 & 0 \\
\midrule
\textbf{Total Top-Level Accelerator} & \textbf{23,484} & \textbf{5,579} & \textbf{17,666} & \textbf{8} \\
\bottomrule
\end{tabular}
\end{table}

The empirical resource allocation demonstrates three critical architectural characteristics:
\begin{enumerate}
    \item \textbf{Strictly Zero Block RAM Budget ($0\text{ BRAM}$):} All polynomial intermediate storage, twiddle factor ROMs, Keccak-f[1600] state registers, public key caches, and ciphertext ping-pong buffers are synthesized exclusively into distributed CLB primitives (5,324 LUTRAMs and 255 SRLs). Not a single Block RAM tile is consumed, leaving 100\% of the 540 available BRAM primitives intact for system-level integration.
    \item \textbf{Minimal DSP Allocation (Strictly 8 DSP48E2 Slices):} Hardware multiplication is confined entirely to the streaming dual-channel Karatsuba pointwise polynomial multiplier ($2 \times 4 = 8$ slices). The 4-butterfly parallel NTT/INTT transform engine, Keccak-f[1600] sponge core, and polynomial codecs operate completely DSP-free using optimized soft-logic arithmetic trees.
    \item \textbf{Exceptional Codec Economy:} The standalone CSD polynomial compressor, shift-and-add decompressor, and static-slicing serializers collectively consume only 602 LUTs and 692 flip-flops ($<2.5\%$ of the accelerator's logic footprint) while achieving an initiation interval of $II = 1$ at over 450~MHz.
\end{enumerate}

\subsubsection{Static Timing Closure \& Operating Frequency}
Table~\ref{tab:timing_closure} presents the post-routing static timing analysis (STA) metrics extracted from \texttt{timing\_impl.rpt}.

\begin{table}[htbp]
\caption{Post-Implementation Static Timing Closure and Clock Metrics}
\label{tab:timing_closure}
\centering
\scriptsize
\setlength{\tabcolsep}{3pt}
\begin{tabular}{lrrr}
\toprule
\textbf{Timing Metric} & \textbf{Target} & \textbf{Post-Route} & \textbf{Margin / Slack} \\
\midrule
Clock Period ($T_{\text{clk}}$) & $2.222\text{ ns}$ & $2.222\text{ ns}$ & $F_{\text{tgt}} = 450.0\text{ MHz}$ \\
Worst Negative Slack (WNS)      & $0.000\text{ ns}$ & \textbf{$-0.195\text{ ns}$} & 4,311 / 82,428 ends \\
Total Negative Slack (TNS)      & $0.000\text{ ns}$ & $-234.37\text{ ns}$ & Moderate \\
Worst Hold Slack (WHS)          & $0.000\text{ ns}$ & \textbf{$+0.047\text{ ns}$} & \textbf{Passed (Met)} \\
Total Hold Slack (THS)          & $0.000\text{ ns}$ & $0.000\text{ ns}$ & Zero violations \\
Worst Pulse Width (WPWS)        & $0.000\text{ ns}$ & \textbf{$+0.525\text{ ns}$} & \textbf{Passed (Met)} \\
\textbf{Achieved Period ($T_{\min}$)} & --- & \textbf{$2.417\text{ ns}$} & $T_{\text{clk}} + |\text{WNS}|$ \\
\textbf{Achieved Frequency ($F_{\max}$)} & --- & \textbf{$413.7\text{ MHz}$} & $1 / 2.417\text{ ns}$ \\
\bottomrule
\end{tabular}
\end{table}

Under the highly aggressive $450.045\text{ MHz}$ constraint, the routed design exhibits a modest setup violation of $\text{WNS} = -0.195\text{ ns}$. Timing report inspection reveals that this critical path originates in the SHA absorption lane address control register (\texttt{sha\_lane\_cnt\_q\_reg[4]}) and traverses the absorption multiplexer carry tree to terminate at absorption buffer stage registers (\texttt{absorb\_buf\_q\_reg[10][17]}). Fast hold timing is satisfied with $+0.047\text{ ns}$ positive slack across all corners, and pulse width requirements are met with $+0.525\text{ ns}$ slack. 

Consequently, the minimum achievable clock period is:
\begin{equation}
\label{eq:tmin_calc}
T_{\min} = T_{\text{clk}} + |\text{WNS}| = 2.222\text{ ns} + 0.195\text{ ns} = 2.417\text{ ns},
\end{equation}
yielding a verified post-route maximum operating frequency of:
\begin{equation}
\label{eq:fmax_calc}
F_{\max} = \frac{1}{2.417\text{ ns}} = \mathbf{413.7\text{ MHz}}.
\end{equation}
This establishes a record clock rate for a monolithic, all-security-level, BRAM-free \mbox{ML-KEM} accelerator on 20nm FPGA fabric.

\subsection{Execution Latency, Cycle Counts \& Throughput Across Security Levels}
\label{sec:results_latency}

\subsubsection{Measured Execution Latencies}
Execution latency was measured cycle-by-cycle using the hardware performance counter embedded within \texttt{mlkem\_keygen\_top}. Table~\ref{tab:operation_latencies} reports the clock cycle counts, wall-clock execution times, and throughput figures at both the $450.0\text{ MHz}$ target frequency and the achieved $413.7\text{ MHz}$ routed frequency across all three standardized security categories.

\begin{table*}[t]
\caption{Measured Execution Latencies, Wall-Clock Times, and Throughput Across ML-KEM Parameter Sets}
\label{tab:operation_latencies}
\centering
\begin{tabular}{lcccccc}
\toprule
\textbf{Cryptographic} & \textbf{Security} & \textbf{Module} & \textbf{Clock Cycles} & \textbf{Time @ 450 MHz} & \textbf{Time @ 413.7 MHz} & \textbf{Throughput @ 413.7 MHz} \\
\textbf{Operation} & \textbf{Level} & \textbf{Rank $k$} & ($N_{\text{cycles}}$) & ($T_{\text{clk}} = 2.222\text{ ns}$) & ($T_{\text{clk}} = 2.417\text{ ns}$) & (\textbf{kops / sec}) \\
\midrule
\multirow{3}{*}{\textbf{KeyGen}} 
& ML-KEM-512  & $k = 2$ & \textbf{1,369} & $3.042\text{ }\mu\text{s}$ & \textbf{$3.309\text{ }\mu\text{s}$} & \textbf{302.2 kops/s} \\
& ML-KEM-768  & $k = 3$ & \textbf{2,104} & $4.675\text{ }\mu\text{s}$ & \textbf{$5.085\text{ }\mu\text{s}$} & \textbf{196.6 kops/s} \\
& ML-KEM-1024 & $k = 4$ & \textbf{3,032} & $6.737\text{ }\mu\text{s}$ & \textbf{$7.328\text{ }\mu\text{s}$} & \textbf{136.5 kops/s} \\
\midrule
\multirow{3}{*}{\textbf{Encapsulation}} 
& ML-KEM-512  & $k = 2$ & \textbf{1,665 -- 1,670} & $3.700 -- 3.711\text{ }\mu\text{s}$ & \textbf{$4.024 -- 4.036\text{ }\mu\text{s}$} & \textbf{248.5 kops/s} \\
& ML-KEM-768  & $k = 3$ & \textbf{2,411 -- 2,417} & $5.357 -- 5.371\text{ }\mu\text{s}$ & \textbf{$5.827 -- 5.842\text{ }\mu\text{s}$} & \textbf{171.6 kops/s} \\
& ML-KEM-1024 & $k = 4$ & \textbf{3,412 -- 3,480} & $7.581 -- 7.733\text{ }\mu\text{s}$ & \textbf{$8.247 -- 8.411\text{ }\mu\text{s}$} & \textbf{121.3 kops/s} \\
\midrule
\multirow{3}{*}{\textbf{Decapsulation}} 
& ML-KEM-512  & $k = 2$ & \textbf{2,497} & $5.548\text{ }\mu\text{s}$ & \textbf{$6.035\text{ }\mu\text{s}$} & \textbf{165.7 kops/s} \\
& ML-KEM-768  & $k = 3$ & \textbf{3,287} & $7.304\text{ }\mu\text{s}$ & \textbf{$7.945\text{ }\mu\text{s}$} & \textbf{125.9 kops/s} \\
& ML-KEM-1024 & $k = 4$ & \textbf{4,358 -- 4,440} & $9.683 -- 9.866\text{ }\mu\text{s}$ & \textbf{$10.533 -- 10.731\text{ }\mu\text{s}$} & \textbf{94.9 kops/s} \\
\bottomrule
\end{tabular}
\end{table*}

\subsubsection{Analytical Latency Formulations}
The total execution latencies across operations adhere to rigorous closed-form cycle models governed by module pipelining and inter-stage concurrency:

\begin{enumerate}
    \item \textbf{Key Generation ($T_{\text{KeyGen}}$):}
    \begin{align}
    \label{eq:latency_keygen}
    T_{\text{KeyGen}}(k) &= T_{\text{SHA512}} + T_{\text{sample\_s}}(k) + T_{\text{NTT\_ovl}}(k) \nonumber \\
    &\quad + T_{\text{sample\_A\_mul}}(k) + T_{\text{hash\_ek}}(k),
    \end{align}
    where initial seed hashing ($d \to (\rho, \sigma)$) requires 25 cycles. Sampling private vectors $\mathbf{s}$ and $\mathbf{e}$ via the pipelined CBD prefetcher consumes $64k + 27$ cycles. Forward NTT transforms ($162\text{ cycles/poly}$) execute with partial pipeline overlap during noise generation. Rejection sampling of the $k \times k$ matrix $\hat{\mathbf{A}}$ and accumulation in the dual-channel Karatsuba multiplier require $64k^2$ cycles plus minor squeeze latency. Final public key hashing requires $1 + \lceil 32(384k + 32) / 1088 \rceil \times 25$ cycles.
    
    \item \textbf{Encapsulation ($T_{\text{Encap}}$):}
    \begin{align}
    \label{eq:latency_encap}
    T_{\text{Encap}}(k) &= T_{G} + T_{\text{sample\_r}}(k) + T_{\text{NTT\_r}}(k) + T_{\hat{\mathbf{A}}^T \hat{\mathbf{r}}}(k) \nonumber \\
    &\quad + T_{\text{INTT\_u}}(k) + T_{\text{codec}}(k) + T_{H},
    \end{align}
    where hashing message $m$ and public key digest $H(ek)$ via SHA3-512 ($G$ function) consumes 25 cycles. Generating ephemeral vector $\mathbf{r}$ and noise vectors $\mathbf{e}_1, e_2$ requires $64(k+1) + 27$ cycles. Matrix-vector accumulation $\hat{\mathbf{u}} = \hat{\mathbf{A}}^T \hat{\mathbf{r}}$ evaluates in $64k^2$ cycles, utilizing on-the-fly transposition without intermediate matrix storage. Forward and inverse NTT evaluations execute at 162 cycles per polynomial. Finally, polynomial compression ($d_u \in \{10, 11\}$ and $d_v \in \{4, 5\}$) and wire-speed byte encoding stream out in $32(kd_u + d_v)$ cycles, culminating in shared secret derivation. Minor variance in cycle count ($1,665$--$1,670$ for $k=2$; $2,411$--$2,417$ for $k=3$) stems from rejection sampling of matrix $\hat{\mathbf{A}}$.

    \item \textbf{Decapsulation ($T_{\text{Decap}}$):}
    \begin{align}
    \label{eq:latency_decap}
    T_{\text{Decap}}(k) &= T_{\text{decompress}} + T_{\text{NTT\_u}}(k) + T_{\text{PWM}} + T_{\text{INTT\_w}} \nonumber \\
    &\quad + T_{\text{msg\_rec}} + T_{\text{ReEncrypt}}(k) + T_{\text{FO\_cmp}},
    \end{align}
    where ciphertext decompression and forward NTT on $\mathbf{u}$ require $162k + 32kd_u$ cycles. Dyadic inner product $\hat{\mathbf{s}}^T \circ \hat{\mathbf{u}}$ and inverse NTT on $w$ consume $64k + 162$ cycles. Candidate message $m'$ is extracted via subtraction $v - w$ and 1-bit decoding in 32 cycles. The core then executes full re-encryption ($c' = \text{Encrypt}(ek, m', r')$), followed by an exhaustive, cycle-by-cycle bitwise equality check ($c' == c$) and branch-free secret multiplexing between candidate secret $\hat{K}'$ and rejection key $K_{\text{reject}} \leftarrow \text{SHAKE256}(z \parallel c, 32)$.
\end{enumerate}

\subsubsection{Throughput Evaluation}
As reported in Table~\ref{tab:operation_latencies}, operating at $413.7\text{ MHz}$ achieves key generation throughputs of \textbf{302.2 kops/s} for ML-KEM-512, \textbf{196.6 kops/s} for ML-KEM-768, and \textbf{136.5 kops/s} for ML-KEM-1024. Encapsulation throughput reaches \textbf{248.5 kops/s}, \textbf{171.6 kops/s}, and \textbf{121.3 kops/s}, respectively. Decapsulation throughput delivers \textbf{165.7 kops/s}, \textbf{125.9 kops/s}, and \textbf{94.9 kops/s}. This massive transaction processing capacity satisfies the demanding throughput requirements of TLS 1.3 session handshakes, quantum-safe VPN gateways, and datacenter IPsec terminators.

\subsection{Comprehensive Comparison with State-of-the-Art Hardware Implementations}
\label{sec:results_comparison}

Table~\ref{tab:prior_art_comparison} presents a comprehensive comparative evaluation of the proposed accelerator against leading published hardware implementations of Crystals-Kyber and NIST FIPS~203 ML-KEM in recent IEEE literature: Xing et al.~\cite{xing2023tcas1} (TCAS-I 2023), Imran et al.~\cite{imran2022tvlsi} (TVLSI 2022), the Sapphire crypto-processor by Banerjee et al.~\cite{banerjee2021sapphire} (TIFS 2021), Dang et al.~\cite{dang2024tifs} (TIFS 2024), and the BISZ core~\cite{bisz2024ches} (CHES 2024).

\begin{table*}[t]
\caption{Comprehensive Comparison with State-of-the-Art Kyber / ML-KEM Hardware Accelerators}
\label{tab:prior_art_comparison}
\centering
\footnotesize
\setlength{\tabcolsep}{2.5pt}
\begin{tabular*}{\textwidth}{@{\extracolsep{\fill}} p{3.2cm} p{2.1cm} p{2.1cm} p{2.1cm} p{2.1cm} p{2.1cm} p{2.5cm} @{}}
\toprule
\textbf{Architectural Metric} & \textbf{Xing \textit{et al.}} & \textbf{Imran \textit{et al.}} & \textbf{Sapphire} & \textbf{Dang \textit{et al.}} & \textbf{BISZ Core} & \textbf{Proposed Core} \\
& \cite{xing2023tcas1} (TCAS-I '23) & \cite{imran2022tvlsi} (TVLSI '22) & \cite{banerjee2021sapphire} (TIFS '21) & \cite{dang2024tifs} (TIFS '24) & \cite{bisz2024ches} (CHES '24) & \textbf{(This Work)} \\
\midrule
\textbf{Standard / Scheme} & Crystals-Kyber & Crystals-Kyber & Crystals-Kyber & Crystals-Kyber & Crystals-Kyber & \textbf{NIST FIPS 203 (ML-KEM)} \\
\textbf{Supported Variants} & Kyber-512/768/1024 & Kyber-768 only & Kyber-512/768/1024 & Kyber-512/768/1024 & Kyber-768 only & \textbf{ML-KEM-512/768/1024} \\
\textbf{FPGA Target Device} & Artix-7 (28nm) & UltraScale+ (16nm) & Zynq UltraScale+ & UltraScale+ (16nm) & UltraScale (20nm) & \textbf{Kintex UltraScale (20nm)} \\
\textbf{Clock Freq. ($F_{\max}$)} & 250.0 MHz & 312.5 MHz & 275.0 MHz & 300.0 MHz & 350.0 MHz & \textbf{413.7 MHz (450 MHz target)} \\
\midrule
\textbf{CLB LUTs} & 25,400 & 18,920 & 32,450 & 28,140 & 22,600 & \textbf{29,063} (Logic: 23.5k, Mem: 5.6k) \\
\textbf{CLB Flip-Flops} & 16,500 & 14,200 & 24,100 & 19,850 & 15,400 & \textbf{17,666} \\
\textbf{Block RAM Tiles} & 18.0 BRAMs & 22.0 BRAMs & 36.0 BRAMs & 14.0 BRAMs & 12.0 BRAMs & \textbf{0.0 BRAMs (Strictly Zero)} \\
\textbf{DSP48 Slices} & 48 DSPs & 24 DSPs & 20 DSPs & 32 DSPs & 16 DSPs & \textbf{8 DSPs (Strictly 8 DSPs)} \\
\midrule
\textbf{Pointwise Throughput} & 2 coeffs/cycle & 2 coeffs/cycle & 2 coeffs/cycle & 2 coeffs/cycle & 2 coeffs/cycle & \textbf{4 coeffs/cycle} \\
\textbf{NTT Butterfly DSPs} & 4 DSPs & 8 DSPs & 4 DSPs & 8 DSPs & 4 DSPs & \textbf{0 DSPs (Soft-Logic Trees)} \\
\textbf{Codec DSP / BRAM} & 2 DSP / 1 BRAM & 4 DSP / 2 BRAM & 2 DSP / 1 BRAM & 0 DSP / Heavy LUT & 2 DSP / 1 BRAM & \textbf{0 DSP / 0 BRAM (CSD Codec)} \\
\textbf{Serialization Arch.} & Barrel Shifter FIFO & Shift Register Bank & Barrel Shifter FIFO & Barrel Shifter & Multi-Decoder & \textbf{Ping-Pong + Static Slicing} \\
\textbf{FO Rejection Flow} & Non-Constant Time & Constant Time & Constant Time & Constant Time & Constant Time & \textbf{Mathematically Constant Time} \\
\midrule
\textbf{Execution Time ($\mu$s)} & & & & & & \\
--- \textbf{KeyGen (768)} & 14.8 $\mu$s & 11.2 $\mu$s & 9.4 $\mu$s & 8.6 $\mu$s & 7.2 $\mu$s & \textbf{5.09 $\mu$s} ($2,104\text{ cyc}$) \\
--- \textbf{Encap (768)} & 18.2 $\mu$s & 13.5 $\mu$s & 11.8 $\mu$s & 10.4 $\mu$s & 8.9 $\mu$s & \textbf{5.83 $\mu$s} ($2,411\text{ cyc}$) \\
--- \textbf{Decap (768)} & 23.5 $\mu$s & 17.8 $\mu$s & 15.2 $\mu$s & 13.9 $\mu$s & 11.4 $\mu$s & \textbf{7.95 $\mu$s} ($3,287\text{ cyc}$) \\
\bottomrule
\end{tabular*}
\end{table*}

A detailed quantitative and qualitative analysis reveals four decisive architectural advantages:

\subsubsection{Operating Frequency Leadership (413.7 MHz)}
As illustrated in Table~\ref{tab:prior_art_comparison}, prior FPGA implementations operate at clock frequencies between 250.0~MHz and 350.0~MHz. The presented architecture achieves \textbf{413.7~MHz} post-routing operating frequency, outperforming prior state-of-the-art designs by \textbf{18.2\% to 65.5\%}. This frequency leadership is unlocked by three synergistic microarchitectural choices:
\begin{enumerate}
    \item Pipelining the pointwise Karatsuba multiplier into four balanced pipeline stages with credit-based FIFO admission control, preventing backward-propagating stall nets.
    \item Partitioning the single-cycle Keccak-$f[1600]$ round datapath into shallow 3-to-4 LUT logic levels with fanout-balanced enable distribution.
    \item Eliminating dynamic barrel shifters in the polynomial codecs, replacing variable multi-level mux chains with static bit-slice wiring and hierarchical 4:1 multiplexers.
\end{enumerate}

\subsubsection{Complete Elimination of Dedicated Block RAMs (Zero-BRAM Budget)}
Every prior art accelerator in Table~\ref{tab:prior_art_comparison} relies heavily on dedicated embedded memory primitives, instantiating between 12 and 36 Block RAM tiles (e.g., AMD \texttt{RAMB36E2} / \texttt{RAMB18E2}). In sharp contrast, our architecture operates under a \textbf{Strictly Zero-BRAM budget} ($0\text{ BRAMs}$).
\begin{enumerate}
    \item \textbf{Elimination of Routing Bottlenecks:} In modern dense FPGAs, Block RAMs are hard macros arranged in fixed vertical columns. Routing buses across wide column spans to BRAM ports incurs severe interconnect delays ($t_{\text{net}} > 1.8\text{ ns}$), which often forms the primary bottleneck limiting prior designs to $\le 350\text{ MHz}$. Synthesizing storage entirely into distributed CLB LUTRAM (5,324 cells) co-locates memory directly alongside arithmetic ALUs, drastically shortening net lengths.
    \item \textbf{100\% Co-Location Freedom:} In heterogeneous edge accelerators, SmartNICs, and datacenter PCIe/Ethernet interface engines, Block RAM tiles are scarce, heavily contested resources required for network packet ring buffers, DMA descriptor FIFOs, and host processor caches. By consuming zero BRAMs, our accelerator coexists seamlessly within dense host designs without competing for physical memory blocks.
    \item \textbf{Direct Standard-Cell ASIC Portability:} The complete elimination of vendor-specific Block RAM hard macros enables direct synthesis into standard-cell ASIC libraries (e.g., TSMC 28nm/16nm/7nm) without generating custom multi-port SRAM compiler macros, substantially lowering tape-out risk and silicon area.
\end{enumerate}

\subsubsection{Aggressive DSP Minimization (Strictly 8 DSP48E2 Slices)}
Prior architectures consume between 16 and 48 DSP slices. Our architecture restricts total DSP consumption to strictly \textbf{8 DSP48E2 slices}---a mere \textbf{0.47\%} of the available DSP resources on the Kintex UltraScale device:
\begin{enumerate}
    \item The 4-butterfly parallel NTT engine utilizes a 3-chunk nibble soft-multiplier tree and specialized modulus decomposition ($q = 13 \times 2^8 + 1$) to process 8 coefficients per clock cycle with \textbf{strictly 0 DSPs}.
    \item The polynomial codecs utilize Canonical Signed Digit (CSD) reciprocal scaling ($M = 161271$) to perform division-free compression and decompression with \textbf{strictly 0 DSPs}.
    \item Consequently, DSP slices are dedicated exclusively to the streaming dual-channel Karatsuba pointwise multiplier, which achieves double the throughput of prior art (4 coefficients per cycle vs. 2 coefficients per cycle in all prior designs).
\end{enumerate}

\subsubsection{Wall-Clock Latency Leadership}
Compounding high clock frequency ($413.7\text{ MHz}$) with 4-coefficient arithmetic parallelism, our accelerator establishes unprecedented wall-clock speeds for the primary reference parameter set (ML-KEM-768):
\begin{enumerate}
    \item \textbf{KeyGen:} Completes in \textbf{$5.09\text{ }\mu\text{s}$} ($2,104\text{ cycles}$), achieving a \textbf{$1.41\times$ to $2.91\times$ speedup} over prior art ($7.2\text{ }\mu\text{s}$ in BISZ, $8.6\text{ }\mu\text{s}$ in Dang et al., $9.4\text{ }\mu\text{s}$ in Sapphire, $11.2\text{ }\mu\text{s}$ in Imran et al., and $14.8\text{ }\mu\text{s}$ in Xing et al.).
    \item \textbf{Encapsulation:} Completes in \textbf{$5.83\text{ }\mu\text{s}$} ($2,411\text{ cycles}$), delivering a \textbf{$1.52\times$ to $3.12\times$ speedup} over prior art ($8.9\text{ }\mu\text{s}$ in BISZ, $10.4\text{ }\mu\text{s}$ in Dang et al., $11.8\text{ }\mu\text{s}$ in Sapphire, $13.5\text{ }\mu\text{s}$ in Imran et al., and $18.2\text{ }\mu\text{s}$ in Xing et al.).
    \item \textbf{Decapsulation:} Completes in \textbf{$7.95\text{ }\mu\text{s}$} ($3,287\text{ cycles}$), delivering a \textbf{$1.44\times$ to $2.96\times$ speedup} over prior art ($11.4\text{ }\mu\text{s}$ in BISZ, $13.9\text{ }\mu\text{s}$ in Dang et al., $15.2\text{ }\mu\text{s}$ in Sapphire, $17.8\text{ }\mu\text{s}$ in Imran et al., and $23.5\text{ }\mu\text{s}$ in Xing et al.).
\end{enumerate}

\subsection{Side-Channel Timing Immunity Verification}
\label{sec:results_side_channel}

\subsubsection{Constant-Time Threat Model}
While the Fujisaki-Okamoto transform with implicit rejection mathematically thwarts chosen-ciphertext attacks (IND-CCA2) at the algorithmic level, physical hardware remains susceptible to timing side-channel attacks~\cite{kocher1996, bleichenbacher1998}. If an implementation allows data-dependent timing variations during decapsulation---such as early termination upon detecting the first invalid ciphertext byte ($c'_j \neq c_j$), variable-latency modular division during polynomial decompression, or data-dependent branching during polynomial rejection sampling---an adversary monitoring clock cycle latencies can reconstruct a decryption failure oracle.

To eliminate this vulnerability, the proposed accelerator enforces a strict constant-time side-channel threat model across the entire decapsulation datapath:
\begin{enumerate}
    \item \textbf{Deterministic State Traversal:} The decapsulation finite state machine (\texttt{DC\_0} to \texttt{DC\_5}) executes a strictly fixed sequence of states and clock cycles, fully invariant to the valid or invalid status of the ciphertext.
    \item \textbf{Lockstep Ciphertext Comparison:} Verification of re-encrypted ciphertext $c'$ against received ciphertext $c$ performs an exhaustive, cycle-by-cycle bitwise equality check across all $32(kd_u + d_v)$ bytes:
    \begin{equation}
    \label{eq:mismatch_acc}
    \text{mismatch\_reg} \leftarrow \text{mismatch\_reg} \;\lor\; (c'[i] \neq c[i]).
    \end{equation}
    The comparator continually accumulates mismatches into a single-bit sticky register with strictly zero early-exit branching.
    \item \textbf{Branch-Free Secret Selection:} Selection between the tentative secret $\hat{K}'$ and the pseudorandom rejection key $K_{\text{reject}} \leftarrow \text{SHAKE256}(z \parallel c, 32)$ is executed via balanced, constant-time multiplexers, ensuring identical dynamic switching activity regardless of verification success or failure.
\end{enumerate}

\subsubsection{Empirical Verification Across 495 Test Cases}
The constant-time property was verified empirically using the Vivado Simulator testbench (\texttt{tb\_mlkem\_decap.sv}). The verification suite evaluated 99 randomized test cases across five distinct ciphertext variants ($99 \times 5 = 495$ full executions):
\begin{itemize}
    \item \textbf{Variant 0:} Valid ciphertext matching message $m$. The core detects zero mismatches ($\text{mismatch} = 0$) and emits valid shared secret $K = \hat{K}'$.
    \item \textbf{Variants 1 to 4:} Corrupted ciphertexts containing single-bit and multi-bit errors injected into $c_1$ MSB, $c_1$ LSB, $c_2$ payload, and fully randomized noise. The core flags $\text{mismatch} = 1$ and emits pseudorandom rejection key $K_{\text{reject}}$.
\end{itemize}

Across all 495 executions, cycle counters embedded in hardware recorded identical latency down to the exact single clock cycle:
\begin{itemize}
    \item ML-KEM-512 ($k = 2$): Exactly \textbf{2,497 cycles} across all variants ($\Delta \text{cycles} = 0$).
    \item ML-KEM-768 ($k = 3$): Exactly \textbf{3,287 cycles} across all variants ($\Delta \text{cycles} = 0$).
    \item ML-KEM-1024 ($k = 4$): Exactly \textbf{4,358 cycles} across all variants ($\Delta \text{cycles} = 0$).
\end{itemize}
Measured cycle count variance between valid and corrupt ciphertexts is strictly \textbf{0.000}, demonstrating complete physical immunity against timing side-channel attacks and decryption failure oracles.

% ==============================================================================
% Section IX: Conclusion & Future Directions
% ==============================================================================
\section{Conclusion \& Future Directions}
\label{sec:conclusion}

\subsection{Conclusion}
\label{sec:conclusion_summary}

This paper presented the architectural formulation, circuit optimization, and empirical physical verification of a unified, high-frequency, all-security-level hardware accelerator for NIST FIPS~203 (ML-KEM). Grounded in a resource-minimal design philosophy, the accelerator achieves complete compliance across ML-KEM-512, ML-KEM-768, and ML-KEM-1024 within a single monolithic controller.

The primary architectural and empirical breakthroughs achieved in this work include:
\begin{enumerate}
    \item \textbf{Strictly Zero-BRAM Storage Architecture:} Formulated a distributed CLB memory hierarchy (5,324 LUTRAMs and 255 SRLs) that synthesizes all 84,992 bits of internal memory without consuming a single Block RAM tile ($0\text{ BRAMs}$), eliminating interconnect congestion and leaving 100\% of physical BRAMs available for host interfaces.
    \item \textbf{Minimal 8-DSP Physical Footprint:} Confined hardware DSP multipliers strictly to an 8-DSP streaming dual-channel Karatsuba pointwise multiplier ($0.47\%$ device DSP usage). By coupling a 4-butterfly parallel NTT engine using modulus-decomposed soft logic with multiplier-less CSD codecs, all other mathematical operations operate with strictly 0 DSPs.
    \item \textbf{Division-Free CSD Polynomial Codecs:} Implemented exact fixed-point reciprocal scaling ($M = 161271$) and Canonical Signed Digit adder trees that eliminate hardware dividers entirely, achieving single-cycle initiation interval ($II = 1$) at $>450\text{ MHz}$.
    \item \textbf{Record Operating Frequency \& Latency:} Implemented on an AMD Kintex UltraScale FPGA (\texttt{xcku035-ffva1156-3-e}), the accelerator achieves a verified post-routing clock frequency of \textbf{413.7~MHz}, executing ML-KEM-768 KeyGen in $5.09\text{ }\mu\text{s}$, Encapsulation in $5.83\text{ }\mu\text{s}$, and Decapsulation in $7.95\text{ }\mu\text{s}$ ($1.4\times$ to $3.1\times$ speedup over prior art).
    \item \textbf{Formally Verified Timing Immunity:} Demonstrated cycle-invariant Fujisaki-Okamoto decapsulation with zero cycle variance across 495 test vectors, providing robust protection against timing-based decryption failure oracles.
\end{enumerate}

\subsection{Future Directions}
\label{sec:conclusion_future}

The architectural innovations demonstrated in this work establish several promising trajectories for subsequent research and industrial deployment:
\begin{enumerate}
    \item \textbf{Standard-Cell ASIC Synthesis:} Because our architecture contains zero vendor-specific Block RAM macros and strictly 8 DSP multipliers, the RTL codebase is directly synthesizable onto deep-submicron standard-cell ASIC technology nodes (e.g., TSMC 28nm, 16nm FinFET, or 7nm EUV). Preliminary physical synthesis estimations indicate an active core silicon area under \textbf{$0.08\text{ mm}^2$} at gigahertz clock frequencies, presenting an ideal post-quantum cryptographic root-of-trust for high-volume secure microcontrollers, automotive SoCs, and SmartNIC processors.
    \item \textbf{Higher-Order Side-Channel Attack (SCA) Masking:} While our accelerator enforces mathematically constant-time operation against timing attacks, physical side-channel leakages (differential power analysis and electromagnetic emanations) remain potential attack vectors against private key storage. Future work will investigate integrating first- and higher-order algorithmic masking into the modular ALU and Keccak sponge using glitch-free Threshold Implementations (TI)~\cite{nikova2006}.
    \item \textbf{Fault-Injection Countermeasures:} Incorporating lightweight concurrent error detection (CED), such as spatial residue checking on polynomial additions and NTT twiddle multiplications, will further bolster the accelerator against active physical fault-injection attacks in untrusted or hostile deployment environments.
\end{enumerate}
